\begin{answers}{ch:pytorch}{道具の準備：Python・NumPy・PyTorch のテンソル}
\item $(4,)$、$(3,1)$、$(2,1,1)$ の三つは形 $(2,3,4)$ になり、$(3,)$ と $(2,4)$ はエラーです（実行して確かめました）。
規則のとおり、右端をそろえて比べます。
$(4,)$ は右端が「4 と 4」で同じで、左の二つの軸は長さ 1 とみなされて引き伸ばされます。
$(3,1)$ は「4 と 1」「3 と 3」、$(2,1,1)$ は「4 と 1」「3 と 1」「2 と 2」で、どれも規則 (1)〜(3) で済みます。
$(3,)$ は右端が「4 と 3」で、どちらも 1 でないのでエラーです。
$(2,4)$ は右端の「4 と 4」は通りますが、その左が「3 と 2」でエラーになります。
\item \code{C.sum(dim=1)} は形 $(3,)$ の $(10,4,8)$ で、右端でそろうため、\textbf{各列}がこの数で割られます。
実際、1 行目は $(0.8, 0.25, 0.125)$ になり、各行の和は $(1.175, 0.7, 1.35)$ で 1 になりません（実測）。
正しくは \code{C / C.sum(dim=1, keepdim=True)} とします。
割る数の形が $(3,1)$ になり、各\textbf{行}がその行の合計で割られます。
各行の和はすべて 1 になります。
\item 勾配は $(3,12,27)$ です。
$\partial f/\partial x_i = 3x_i^{2}$ なので、$\vx=(1,2,3)$ を代入すると $(3,12,27)$ で、自動微分の結果と一致します（$f=36$）。
0 に戻さずに $f$ をもう一度計算して \code{backward()} を呼ぶと、前回の分が足し込まれて $(6,24,54)$ になります（実測）。
なお、$f$ を計算し直さずに \code{f.backward()} だけを 2 回呼ぶと、エラーになります。
1 回目の \code{backward()} のあと、計算の記録が捨てられるからです。
\item 自作の \code{softmax} は \code{tensor([[nan, nan, nan]])} を返します（実測）。
\code{float32} が表せる最大の数は約 $3.4\times10^{38}$ で、$e^{89}$ 程度で超えます。
$e^{1000}$ はあふれて \code{inf}（無限大）になり、$\infty/\infty$ が \code{nan}（数でない値）になるのです。
直すには、\code{softmax} の最初の行の前に \code{z = z - z.max(dim=-1, keepdim=True).values} を入れます。
\code{max} は行ごとの最大値（\code{.values}）とその番号を返します。ここでは値だけを使います。
各行から同じ数 $c$ を引いても、$e^{z_i-c}/\sum_j e^{z_j-c}=e^{z_i}/\sum_j e^{z_j}$ なので、答えは変わりません。
最大値を引くと指数の中は 0 以下になり、$e^{z}$ は 1 以下に収まります。
結果は $(0.6652, 0.2447, 0.0900)$ で、\code{torch.softmax} と全桁が一致しました。
\item 勾配は $(-0.3348,\ 0.2447,\ 0.0900)$ で、$\vp$ から「正解の位置だけ 1」のベクトル $(1,0,0)$ を引いたものです（実測）。
softmax の出力は $\vp=(0.6652, 0.2447, 0.0900)$、損失は $-\log 0.6652\approx0.4076$ です。
勾配の 1 番目だけが $0.6652-1$ で、残りは $\vp$ の成分そのものです。
正解の位置は負なので、そのロジットを増やせば損失が下がり、ほかは減らせば下がります。
この関係 $\vp-(\text{正解の位置だけ 1})$ は、\chapref{ch:ml}で式から導きます。
\end{answers}
